\begin{titlepage}
\fontfamily{ptm}\selectfont\centering
\vspace*{-0.6cm}
{\bfseries\large INFORME TÉCNICO — GUÍA–TALLER No.~2\par}
\vspace{0.25cm}
{\bfseries\normalsize Mecanismos de seguridad para redes IEEE 802.3\\ con servidores TCP/IP\par}
\vspace{0.6cm}
\includegraphics[width=7.4cm]{fig/logo-ud-portada.png}\par
\vfill
{\bfseries\itshape\large Juan Esteban Buitrago Chavez - 20211020005\par}
\vfill
{\bfseries\large Docente: OCTAVIO JOSE SALCEDO PARRA\par}
\vfill
{\bfseries\large Facultad de Ingeniería, Universidad Distrital Francisco José de Caldas\par}
\vspace{0.45cm}
{\bfseries\large 020-81 - MECANISMOS DE SEGURIDAD INFORMÁTICA\par}
\vspace{0.45cm}
{\bfseries\large Noviembre de 2026\par}
\vspace{0.8cm}
{\small Repositorio: \url{https://github.com/yatoDx/Laboratorio-seguridad-tcpip}\par}
\end{titlepage}

\chapter*{Resumen}
\thispagestyle{plain}
Este informe desarrolla la Guía–Taller No.~2 de mecanismos de seguridad para redes IEEE~802.3 con servicios TCP/IP. Primero se monta y verifica un laboratorio aislado de cuatro máquinas virtuales Ubuntu Server 24.04.5 LTS (gateway con NTP y firewall nftables, servidor web nginx, DNS dnsmasq y mensajería XMPP Prosody) y un equipo analista Windows con Wireshark, todos en la red 192.168.50.0/24 y con relojes sincronizados.

La Parte~I analiza diez casos críticos: inundación TCP SYN, túnel DNS, exposición de servicios, detección de anomalías con aprendizaje automático, resets TCP, reflexión UDP, ICMP atípico, escaneo de puertos, correlación de alertas al estilo SIEM y seguridad de XMPP sobre TLS. Cada caso tiene su ficha de nueve campos, su evidencia con huella SHA-256, una regla de detección o mitigación y preguntas de análisis. Para este informe se re-ejecutaron los analizadores sobre las mismas capturas y produjeron salidas idénticas byte a byte a las registradas, y el paquete de entrega verifica 242 archivos por copia contra su manifiesto.

La Parte~II implementa un sistema de localización basado en consentimiento: una aplicación web (PWA) con API REST en FastAPI, registro por número telefónico con código de verificación, mapa Leaflet/OpenStreetMap, historial, estadísticas y geocercas. El diseño parte del riesgo de que un localizador se use para acechar: el uso compartido nace apagado, cada observador necesita aprobación expresa con vencimiento, el titular puede revocar al instante y ve quién consultó. La API pasa 61 pruebas automáticas y el flujo completo se probó en Chrome con dos personas.

Los resultados muestran que un indicador aislado no confirma un ataque, que medir sobre la evidencia corrige supuestos (por ejemplo, que 25 SYN no son 25 puertos distintos) y que el aprendizaje automático sirve para priorizar la revisión humana, no para decidir bloqueos.

\vspace{0.6cm}
\noindent\textbf{Index Terms}---seguridad de redes, IEEE 802.3, TCP/IP, Wireshark, análisis de PCAP, nftables, detección de anomalías, aprendizaje automático, SIEM, TLS, XMPP, privacidad, consentimiento.
\par\vspace{0.8cm}
\noindent\fbox{\parbox{\dimexpr\linewidth-2\fboxsep-2\fboxrule\relax}{\small\textbf{Nota:} Este es un trabajo académico desarrollado en equipo personal. Todo ocurrió dentro de un entorno controlado. Cada prueba, paquete y alerta formó parte de un ejercicio académico realizado con equipos propios y máquinas virtuales, utilizando direcciones privadas y tráfico de laboratorio. El trabajo se desarrolló con fines de aprendizaje, sin intervenir sistemas externos ni afectar redes institucionales o de terceros.}}
